% Proposed replacement complexity section. Labels/names may be reconciled by the editor.
\section{Finality, Throughput, and Communication}\label{sec:complexity}
\subsection{Measurement points and steady-state costs}
Let $n_i=|V_i|$, let $c$ be the committee size, and let $m$ be the number
of regions. A new application batch contains at most $\ell$ body bits;
$\kappa$ bounds a signature together with its signer identifier and a
constant number of digests and counters. Certificates below are explicit
signature lists. All message counts mean point-to-point transmissions by
correct participants under the admitted workload; Byzantine wire flooding
is excluded. A batch's mini-block count and body size remain bounded.

Measure a ready batch's latency from its leader's first proposal transmission.
For an honest, already initialized leader in region $i$, define
\[
 C_i=\delta_R+2\Delta_R,
 \qquad F_i=C_i+\delta_R,
 \qquad I_i=2\delta_R.
\]
Here $C_i$ bounds commitment at every correct regional replica, and $F_i$
is a conservative bound for a transferable final receipt at a regional
collector: after the latest honest commitment, one more regional message
delay collects matching post-commit endorsements. Sending the resulting
complete object to proposer $P$ gives $F_i+\delta_P$. Direct endorsements
to $P$ can avoid this extra collection hop. These expressions presume
that bodies, ancestry and verification fit the stipulated complete-message
bounds; otherwise their additional acquisition time must be added.

The leader can obtain a vote QC and propose the next batch after $I_i$,
while earlier commitment timers continue. Thus $I_i$, not $F_i$, is the
ideal batch interval. Construction time, bandwidth and processing can
increase it. If a batch contains $u$ mini-blocks, $u/I_i$ is its sustained
mini-block rate; it does not imply evenly spaced mini-block finalizations
inside a batch. Transaction throughput additionally depends on batch size.

With full proposal forwarding and unaggregated parent QCs, one regional
physical block costs
\[
 \mathsf{Msgs}_i=O(n_i^2),\qquad
 \mathsf{Bits}_i=O(n_i^2\ell+n_i^3\kappa)
\]
messages and bits. The cubic control term arises because an
$O(n_i\kappa)$-bit QC accompanies up to $O(n_i^2)$ forwarded proposals.
Votes and post-commit endorsements contribute $O(n_i^2\kappa)$ additional
bits, already dominated here. Each subscribed external consumer adds
$O(n_i)$ endorsement messages and $O(n_i\kappa)$ control bits per batch;
body retrieval is charged separately. Compact certificates and cached
proposal forwarding may improve these bounds only with their signer,
setup and retrieval costs included.

\subsection{The four operating paths}
Let $X_{ij}$ be the complete-object transfer and verification time for an
already assembled anchor receipt and its missing contents to reach the
correct participants of successor region $j$. It includes any necessary
successor-side dissemination; cached contents reduce its payload. Let
$A_j$ be the time from this common readiness to a current-view certified
anchor. The specified initial NEWVIEW exchange gives the conservative
good-case bound $A_j\leq2\Delta_R+2\delta_R$, once the required readiness
skew is satisfied. A separately specified preinitialization mechanism may
reduce $A_j$; it is not silently assumed to be zero.

Let $D_i$ be any terminal-drain delay beyond the steady-state bound $F_i$.
It is zero on the healthy direct-commit path. Recovery of a terminal batch,
extra failed leaders and missing data are charged to $D_i$ when they occur.
For normal rotation define the proposal-to-proposal boundary interval
\[
 H_{ij}=F_i+D_i+X_{ij}+A_j.
\]
It starts when the old region sends its final application-batch proposal
and ends when the new region sends its first fresh proposal. The first
new transferable receipt is available another $F_j$ later.
Write
\[
 \begin{aligned}
 T_i^{\rm view}&=4\Delta_R+2\delta_R,\\
 T_{ij}^{\rm acc}&=T_{\rm detect}+T_{\rm cut}+T_C+T_{L1}
                   +X_{ij}+A_j+F_j.
 \end{aligned}
\]
The first expression bounds replacement by an honest next leader from the
first correct quit, with state already available. The second lists the
accident phases, defined below.

\begin{table}[H]\centering\small
\begin{tabularx}{\linewidth}{@{}>{\raggedright\arraybackslash}p{.24\linewidth}Y Y@{}}
\toprule
Operating path & Finality or first new final receipt & Batch interval after recovery\\
\midrule
Ideal in-region sequencing & $C_i$ to regional commitment;
$F_i$ to a transferable receipt. & $I_i=2\delta_R$.\\
\addlinespace[3pt]
In-region leader replacement & From the first correct quit:
$T_i^{\rm view}+F_i$. & $I_i$ after the first fresh proposal.\\
\addlinespace[3pt]
Regular term-end rotation $i\to j$ & From the old terminal proposal:
$H_{ij}+F_j$. From its final receipt: $X_{ij}+A_j+F_j$. &
$I_j$ within the new term; one boundary interval $H_{ij}$.\\
\addlinespace[3pt]
Accident rotation $i\to j$ & From the outage:
$T_{ij}^{\rm acc}$.
& $I_j$ after recovery; no rate guarantee while waiting for these premises.\\
\bottomrule
\end{tabularx}
\caption{Finality latency and proposal spacing are different metrics.
Add the final consumer-delivery delay to the receipt endpoints.
The accident expression is a sequential upper bound on named phases,
conditional on each phase completing.}
\label{tab:operating-paths}
\end{table}

The in-region replacement bound follows from the quit protocol. A quit proof
reaches every correct replica within $\Delta_R$; each waits another
$\Delta_R$ before entry; the new leader waits $2\Delta_R$ and then spends
at most $2\delta_R$ re-certifying its anchor. Detection precedes this
measurement. A silent leader's first fresh-proposal deadline is
$6\Delta_R$ after local view entry, but later detection depends on the
deadline index and the attained progress. Additional unsuccessful views
and state transfer are separate costs, not covered by one honest-successor
replacement.

Regular rotation means exhaustion of the fixed quota of $K$ application
batches, not an unsynchronized wall-clock timeout. If a wall-clock tenure
is later added, its authenticated expiry and terminal-cut rules must first
be specified. In a homogeneous fault-free schedule with interval $I$ and
boundary interval $H$, the amortized interval and rate are
\[
 \overline I=\frac{(K-1)I+H}{K},
 \qquad \lambda_{\rm batch}=\frac{1}{\overline I}.
\]
The term has $K-1$ internal proposal gaps and one cross-region gap; adding
the new batch's finality latency again would double-count the pipeline.
For heterogeneous regions, average the corresponding sums of internal
gaps and boundary gaps over the actual schedule. Startup and final drain
are boundary effects for a finite measurement window.

In the accident path, $T_{\rm detect}$ is observation delay,
$T_{\rm cut}$ is safe-cut discovery, old-authority closure and data recovery,
$T_C$ is ordered committee-entry certification, and $T_{L1}$ is the
required L1 inclusion and finality wait. These tasks may overlap; the table
does not claim their sum is a tight lower bound. No finite accident bound
follows without a closed cut, recoverable anchor data, a live successor,
committee progress and L1 progress. A passive status update incurs none
of this handoff delay and does not reset the current quota.

\subsection{Rotation and registry communication}
One in-region replacement adds $O(n_i^2)$ transmissions and
$O(n_i^3\kappa)$ control bits with explicit certificate lists, plus any
missing ancestral data. NEWVIEW re-certification is included; the first
fresh batch then has the steady-state cost above. If terminal drain emits
$d$ additional control blocks, charge their proposal and vote exchanges
separately. They consume no application quota and carry no new application
body, although inherited data may still need transfer.

For one normal handoff, suppose one reachable holder already has the
assembled old receipt and sends it to each successor member. Let $a_{ij}$
be the number of missing anchor/ancestry data bits and $s_{L1}$ the size of
the needed registry-history authentication. The outer transfer costs
\[
 O(n_j)\ \text{messages},\qquad
 O\bigl(n_j(a_{ij}+n_i\kappa+s_{L1})\bigr)\ \text{bits}.
\]
Receipt collection has already been charged to the old batch. Successor
bootstrap adds $O(n_j^2)$ messages and $O(n_j^3\kappa)$ control bits with
the specified all-to-all NEWVIEW procedure. Failed retrieval attempts,
additional holders and uncached history add their actual costs. The
$O(n_j)$ outer transfer is not a claim that the regional consensus itself
has linear communication.

A full status map in the fixed region-list order has $O(m)$ bits; an
entry with an explicit committee signature list has size
$s_E=O(m+c\kappa)$. One status update contributes an $O(s_E)$-bit L1
transaction payload and, if disseminated directly to $N$ interested
participants, $O(N)$ messages and $O(N(s_E+s_{L1}))$ bits. Let $M_C$ and
$B_C$ denote the committee service's internal observation, agreement and
certification costs. They remain explicit parameters because threshold
signatures alone do not specify an agreement protocol. An accident
rotation adds these update costs and its data-recovery traffic to the
normal anchor-transfer and successor-bootstrap costs. A passive update
pays registry costs alone.

\subsection{Rollup verification and storage}
The rollup relation authenticates every included batch, every applicable
registry update and every intervening rotation, in addition to data
correspondence, DA and execution. Its direct verification work therefore
depends on the number of batches, regional signatures, status entries and
transaction bytes in the claimed interval. A validity proof or fraud-proof
implementation determines proving, verification, challenge and on-chain
costs; neither constant L1 cost nor a numerical settlement delay follows
from the relation alone. Sequencer finality, transfer to $P$ and L1
settlement are separate endpoints. Retention grows with the confirmed
history unless an independently justified pruning/checkpoint policy is
installed. All timing bounds require enough bandwidth and processing for
the counted messages and bytes.
